\chapter*{Abstract}
\addcontentsline{toc}{chapter}{Abstract}
Neighbourhood lending ends in a physical handover between people who often have little or no shared history, so a platform has to decide whether a loan is safe from evidence that is sparse, indirect and open to manipulation. Digital Mohalla is a neighbourhood resource-sharing platform built around that decision. Its core is an assurance engine that combines a recency-weighted Beta posterior over a borrower's verified loan outcomes, bounded $k$-hop evidence through the community's trust graph and interpretable manipulation signals, and returns one of four evidence states with its uncertainty, recommended safeguards and a written explanation; uncertain or suspicious cases go to a human moderator instead of being rejected automatically.

Iteration~1 delivered a live core: an event-sourced borrowing lifecycle (14 states; every lifecycle rule, namely allowed transitions, actor roles, no double booking and append-only events, is also enforced in PostgreSQL), the assurance engine, disputes and moderator decisions, a resident and moderator app, and a CI/CD pipeline deploying to a public instance. \NBackendTests{} backend and \NFrontendTests{} frontend tests pass. All seven evaluation scenarios required at mid-term return the expected state and explanation, and eleven of the thirteen scenarios reproduce their expected state end to end on synthetic data. Parameters remain hypotheses until the FYP-2 evaluation, which this report plans.
